% Problem-specific proof note for QBE-OP-OPTCTRL-COLD-CLEAN-001.
% This file is intended for users who want to copy the mathematical
% construction into their own manuscript.

\subsection{A finite block encoding for a transfer operator}

Let \(T\), \(\tau\), and \(S\) be one-qubit registers.  We order the system
basis by
\[
  \operatorname{idx}(T,\tau,S)=4T+2\tau+S ,
  \qquad T,\tau,S\in\{0,1\}.
\]
The target operator is
\[
  E_1
  =
  |0\rangle\!\langle 1|_T
  \otimes
  |0\rangle\!\langle 1|_\tau
  \otimes I_S .
\]
Equivalently, its only nonzero entries are
\[
  E_1[\operatorname{idx}(0,0,0),\operatorname{idx}(1,1,0)] = 1,
  \qquad
  E_1[\operatorname{idx}(0,0,1),\operatorname{idx}(1,1,1)] = 1 .
\]

Add one block-encoding signal qubit \(a\), and order the full basis by
\[
  \operatorname{full}(a,T,\tau,S)
  =
  8a+4T+2\tau+S .
\]
Define a permutation \(\pi\) of the \(16\) full basis states by the table
\[
\begin{array}{c|cccccccccccccccc}
x&0&1&2&3&4&5&6&7&8&9&10&11&12&13&14&15\\
\hline
\pi(x)&8&9&10&11&12&13&0&1&2&3&4&5&6&7&14&15 .
\end{array}
\]
Let \(U_\pi\) be the permutation matrix with
\[
  U_\pi[y,x]=
  \begin{cases}
    1,& y=\pi(x),\\
    0,& \text{otherwise}.
  \end{cases}
\]

\paragraph{Unitarity.}
The table above is a permutation: it has the explicit inverse
\[
\begin{array}{c|cccccccccccccccc}
y&0&1&2&3&4&5&6&7&8&9&10&11&12&13&14&15\\
\hline
\pi^{-1}(y)&6&7&8&9&10&11&12&13&0&1&2&3&4&5&14&15 .
\end{array}
\]
Therefore every row and every column of \(U_\pi\) has exactly one nonzero
entry, so \(U_\pi\) is unitary.

\paragraph{Clean-block calculation.}
The clean block is the block selected by \(a=0\):
\[
  B[row,col] =
  U_\pi[\operatorname{full}(0,row),\operatorname{full}(0,col)] .
\]
For the selected source states,
\[
  \pi(\operatorname{full}(0,1,1,0))=\operatorname{full}(0,0,0,0),
  \qquad
  \pi(\operatorname{full}(0,1,1,1))=\operatorname{full}(0,0,0,1).
\]
Thus the two desired entries of \(B\) equal \(1\).  For every other clean input
column \(\operatorname{full}(0,T,\tau,S)\), the image under \(\pi\) has
signal bit \(a=1\), so it contributes no entry to the \(a=0\) clean block.
Hence \(B=E_1\).

\paragraph{Certified statement.}
The Lean development records this construction as
\[
  \texttt{coldE1CandidateImage},\qquad
  \texttt{coldE1CandidateMatrix}.
\]
The theorem
\[
  \texttt{coldE1Candidate\_blockProjection}
\]
proves that the clean block equals \(E_1\), and
\[
  \texttt{coldE1CandidateImage\_permutation\_certificate}
\]
proves that the candidate image is a permutation.  The certified high-level
resource tuple is
\[
  (\#\text{gates},\text{depth},\#\text{auxiliary qubits},\#\text{oracle calls})
  =
  (4,4,1,0).
\]
The exported Qiskit diagnostic checks the same finite permutation matrix and
passes with zero block-entry and unitarity error.  This diagnostic is a
post-Lean executable check; the reusable mathematical certificate is the Lean
proof above.

